\documentclass[11pt]{article}
\usepackage[a4paper,margin=1in]{geometry}
\usepackage{amsmath,amssymb,mathtools,bm}
\usepackage{enumitem,booktabs,array}
\usepackage{tcolorbox}
\usepackage{hyperref}
\usepackage[protrusion=true,expansion=false]{microtype}
\hypersetup{colorlinks=true,linkcolor=blue,urlcolor=blue,citecolor=blue}
\setlist[itemize]{topsep=2pt,itemsep=2pt}
\setlist[enumerate]{topsep=2pt,itemsep=2pt}
\newtcolorbox{keybox}{colback=blue!3!white,colframe=blue!40!black,boxrule=0.4pt,arc=1mm}
\newtcolorbox{alertbox}{colback=red!3!white,colframe=red!40!black,boxrule=0.4pt,arc=1mm}
\newtcolorbox{answerbox}{colback=green!3!white,colframe=green!40!black,boxrule=0.4pt,arc=1mm}
\newtcolorbox{drillbox}{colback=yellow!5!white,colframe=orange!60!black,boxrule=0.4pt,arc=1mm}
\newcommand{\EV}{\mathbb{E}}
\newcommand{\Var}{\mathrm{Var}}
\newcommand{\Cov}{\mathrm{Cov}}
\newcommand{\blank}[1]{\underline{\hspace{#1}}}

\title{E11-01: Rauh \& Sufi (2010, RFS) \\
\large ``Capital Structure and Debt Structure'' --- Active Recall Workbook}
\author{Empirical Corporate Finance II --- Exam Prep}
\date{\today}
\begin{document}\maketitle

\begin{keybox}
\textbf{Paper in one sentence.} Using hand-collected data on the composition of US corporate debt for a representative sample of firms, RS show that firms use multiple debt \emph{types} (bank loans, senior bonds, subordinated bonds, convertibles, etc.) in predictable patterns tied to firm credit quality, and that accounting for debt structure explains capital-structure puzzles that aggregate leverage studies miss.

\textbf{Placement.} Foundational paper on debt \emph{heterogeneity}, reshapes empirical capital structure research. Widely cited motivator for modern debt-structure theories.
\end{keybox}

\section*{Round 1: Complete Walk-Through}

\subsection*{1.1 Research question}
Most empirical capital-structure papers use book leverage as a single variable. But debt comes in many forms, with different seniority, flexibility, monitoring properties. Do firms choose debt structure systematically, and does ignoring this obscure capital-structure tests?

\subsection*{1.2 Data}
\begin{itemize}
\item Hand-collect 10-K debt-footnotes for 305 firms randomly sampled over 2000--2005.
\item Categorize debt into: commercial paper, revolving credit, term loans, senior bonds, subordinated bonds, convertible bonds, capital leases, other.
\item Merge with Compustat fundamentals, Moody's credit ratings.
\end{itemize}

\subsection*{1.3 Key descriptive facts}
\begin{itemize}
\item 70\% of firms use 2 or more debt types; $\sim$30\% use 3+.
\item Composition correlates strongly with credit quality:
\begin{itemize}
\item AA/A firms: heavy commercial paper, long-term bonds.
\item BBB: bank revolvers + senior bonds.
\item BB/B: term loans + senior + subordinated.
\item CCC: junior and convertible issuance, capital leases.
\end{itemize}
\item Dispersion in debt structure is as large as dispersion in book leverage.
\end{itemize}

\subsection*{1.4 Empirical results}
\begin{itemize}
\item Single-debt-type regressions would miss meaningful cross-sectional variation.
\item Trade-off theory explains book leverage modestly ($R^2 \sim 20\%$); adding debt-type composition boosts explanatory power substantially.
\item Credit-quality gradient in debt types is robust across specifications.
\end{itemize}

\subsection*{1.5 Mechanism}
\begin{itemize}
\item Different debt types offer different combinations of flexibility, monitoring, and cost.
\item High-quality firms optimize for cheap capital (CP, bonds); low-quality firms need monitored or secured debt (bank, leases).
\item Debt structure co-determined with leverage by firm characteristics.
\end{itemize}

\subsection*{1.6 Implications}
\begin{itemize}
\item Capital-structure research should jointly model leverage and debt composition.
\item Ratings agencies implicitly use debt-structure information; rating-based proxies pick this up.
\item Motivates theoretical work on multi-debt-type choice (Diamond 1991, Rauh-Sufi 2014).
\end{itemize}

\subsection*{1.7 Caveats}
\begin{alertbox}
\begin{itemize}
\item Hand-collected sample small relative to Compustat.
\item 10-K footnote data may omit off-balance-sheet obligations.
\item Limited time dimension (5 years).
\end{itemize}
\end{alertbox}

\section*{Round 2: Level 1 Fill-in}
\begin{enumerate}
\item Data: \blank{2cm}-collected 10-K footnotes, \blank{1cm} firms.
\item Fraction using $\geq$2 debt types: \blank{1cm}\%.
\item Credit-rating gradient: high-quality $\to$ \blank{2cm}; low-quality $\to$ secured/junior.
\item Adding composition to capital-structure regression boosts \blank{2cm}.
\item Motivates modeling leverage \emph{and} debt \blank{2cm}.
\end{enumerate}

\section*{Round 3: Level 2 Prompts}
\begin{enumerate}
\item Describe the data and sampling procedure.
\item Report key descriptive facts about debt composition.
\item Discuss the credit-quality gradient.
\item How does debt structure complement leverage in capital-structure regressions?
\item What theoretical directions does RS motivate?
\end{enumerate}

\section*{Round 4: Minimal Scaffolding}
From a blank page: (i) question, (ii) data, (iii) facts, (iv) regression gains, (v) implications.

\section*{Round 5: Blank Page}
Essay: debt heterogeneity and capital structure.

\vspace{6pt}
\begin{drillbox}
\textbf{Speed Drill.} (1) Data source. (2) Sample size. (3) \% using multiple debt types. (4) Credit-quality pattern. (5) Policy implication.
\end{drillbox}

\section*{Quick Reference \& Answer Key}
\begin{answerbox}
\begin{itemize}
\item \textbf{Data.} Hand-coded 10-K footnotes, 305 firms.
\item \textbf{Fact.} 70\% use $\geq$2 debt types.
\item \textbf{Gradient.} CP/bonds $\to$ bank loans $\to$ junior/secured.
\item \textbf{Lesson.} Jointly model leverage \& composition.
\end{itemize}
\end{answerbox}

\begin{keybox}
\textbf{30-second exam pitch.} Rauh \& Sufi (2010) build a hand-collected dataset of corporate debt composition from 10-K footnotes of 305 randomly-sampled US firms 2000--2005, cataloguing commercial paper, revolvers, term loans, senior/subordinated bonds, convertibles, capital leases, and other debt. Two-thirds of firms use multiple debt types, with a clear credit-quality gradient: investment-grade firms rely on commercial paper and long-term bonds, whereas low-rated firms rely on bank loans, subordinated and convertible debt, and capital leases. Aggregating all debt into a single leverage variable misses cross-sectional variation that is economically large. Joint modelling of leverage and debt composition substantially raises explanatory power in capital-structure regressions. The paper reshaped empirical capital-structure research and motivated theoretical work on multi-instrument debt choice, including Diamond-Rajan theories and subsequent Rauh-Sufi follow-ups on debt-structure-based leverage targets.
\end{keybox}

\end{document}
